\section{Repeated-prediction error and dependence}\label{app:prediction}
For a fixed fitted classifier and fixed labels, $\Delta$ is fixed. To discuss
bias and variance, specify a repeated-prediction process, such as repeated
training samples, while fixing $y,w$ and the inputs that determine $w$.
All moments below
condition on this fixed design. Define $b=\E[e]$ and the full joint covariance
$\Sigma=\Cov(e)$, without assuming independence across items. Then
\begin{align}
 \operatorname{Bias}(\widehat\theta)&=w^{\mathsf T}b,\notag\\
 \Var(\Delta)&=w^{\mathsf T}\Sigma w
 =\sum_jw_j^2\Sigma_{jj}
   +2\sum_{j<\ell}w_jw_\ell\Sigma_{j\ell},\label{eq:targetmoments}\\
 \E[\Delta^2]&=(w^{\mathsf T}b)^2+w^{\mathsf T}\Sigma w.
 \label{eq:targetrisk}
\end{align}
Equations~\eqref{eq:estimand}--\eqref{eq:targetrisk} connect the research quantity
to its item weights, systematic error, and uncertainty. Covariance between
item errors enters the estimate's variance. Dependence alone does not determine
bias; systematic errors on items with nonzero weights can make
$w^{\mathsf T}b$ nonzero.

Browsing totals and regression vectors are corresponding joint quantities.
Writing $\widehat T=Cq$, $d=\widehat T-T$, and $\widehat\beta=ACq$ gives
\begin{align}
 d&=Ce,\label{eq:identity}\\
 \E[d]&=Cb, & \Cov(d)&=C\Sigma C^{\mathsf T},\label{eq:moments}\\
 \E[\widehat\beta-\beta^*]&=ACb,
 &\Cov(\widehat\beta-\beta^*)&=AC\Sigma C^{\mathsf T}A^{\mathsf T}.
 \label{eq:regression}
\end{align}
The same item error is reused across all visits and users. These formulas
quantify the specified prediction-error process; they do not include browsing-log
error, user sampling, or uncertainty about the construct. They concern predicted
exposure as an outcome. Noisy exposure as a regressor is a different
errors-in-variables problem treated in Section~\ref{sec:covariates}.

\subsection{Dependence follows the item weights}
There are two sources of dependence. An item's predicted label can be reused
across records, and different items can be classified by the same fitted model.
Shared training data, model parameters, and similar item features can cause
errors to move together under repeated fitting. Distinct items do not establish
independent errors. In browsing, users need not visit the same domains to have
correlated measurement errors: errors on different visited domains can covary.
In equation~\eqref{eq:targetmoments}, positively correlated mistakes increase
variance when their item weights have the same sign, and decrease it when
the weights have opposite signs.

For a limiting example, suppose centered errors are shared within item
families: $e_j-b_j=u_{g(j)}$, with independent family shocks $u_g$ of variance
$\tau_g^2$. This is compatible with binary predictions when the classifier
assigns one common random prediction to every item in a family. Then
\begin{equation}\label{eq:families}
 \Var(\Delta)=\sum_g\tau_g^2
                 \left(\sum_{j:g(j)=g}w_j\right)^2.
\end{equation}
A family of $m$ items with unit weights contributes $m^2\tau_g^2$, whereas
assuming independent item errors would give $m\tau_g^2$. Signed weights can
instead cancel within a family. The relevant concentration is in aggregate
family influence, not just individual item influence. This example is a
special covariance structure, not an assertion that the right families are
known or that all within-family mistakes are perfectly correlated.

Independent item errors are another special case: $\Sigma$ is diagonal,
with entries $v_j$. For browsing totals, equation~\eqref{eq:moments} reduces to
\begin{equation}\label{eq:covariance}
 \Var(d_i)=\sum_jc_{ij}^2v_j,\qquad
 \Cov(d_i,d_\ell)=\sum_jc_{ij}c_{\ell j}v_j.
\end{equation}
Let $\alpha_j$ be the false-positive probability on a truly negative item
and $\beta_j$ the false-negative probability on a truly positive item.
For binary predictions, $v_j=\alpha_j(1-\alpha_j)$ on negative items and
$v_j=\beta_j(1-\beta_j)$ on positive items. In browsing, repeated visits multiply one
classification error; treating visits as independent would incorrectly replace
squared visit weights by visit counts. Even this special case induces dependence
across users who share uncertain domains.

An uncertainty procedure must follow the same structure as the estimand.
Resampling downstream records while holding predictions fixed does not propagate
repeated-training uncertainty. Each replicate must preserve reused labels across
records and the joint distribution of errors across items. Independently
perturbing each item label imposes a diagonal $\Sigma$; resampling the classifier
fit and predicting all items together can preserve dependence induced by that
fitting process, provided the resampling scheme represents that process.
A prediction-error covariance cannot simply be added to an arbitrary reported
regression covariance: the components may overlap or be dependent. Neither
larger downstream samples nor clustering records alone guarantees that this error
disappears,
and standard errors do not repair bias.

\subsection{Many items do not ensure a normal approximation}
Even with independent item errors, diffuse variance shares alone do not
justify Gaussian inference. Let all true labels be zero and let each of $k$
items be falsely labeled positive independently with probability $1/k$.
With $w_j=1$, total error $B_k$ is binomial with mean one and variance $1-1/k$.
The largest item variance share is $1/k$, yet
$B_k-1\Rightarrow\operatorname{Poisson}(1)-1$, not a normal distribution.
At $k=10{,}000$, an oracle interval using the exact bias and variance but a
normal 1.96 cutoff has coverage \RareNormalCoverage{}\%. This calculation
uses exact binomial probabilities, not simulation.

For independent bounded binary errors, a sufficient Lindeberg condition is
\[\frac{\max_j|w_j|}{\sqrt{\sum_jw_j^2v_j}}\to0.\]
A small maximum variance share
$\max_j w_j^2v_j/\sum_jw_j^2v_j$ is weaker when errors become rare. The example
violates the former condition and satisfies the latter. A growing number of
items cannot substitute for a condition on the distribution of their errors.
This is a prediction-error example, distinct from the fixed-label audit design.
